\chapter{Corps}

\section{construction de corps}

\textbf{Le corps des fractions.} Pour cela il faut partir d'un anneau commutatif intègre. Lorsque $A$ est intègre, on peut considérer les fractions d'éléments de $A$
$$\text{Frac}(A) \triangleq \left\{\frac{a}{b} \mid a, b \in A, b \neq 0\right\},$$
avec comme d'habitude $\frac{a}{b} = \frac{c}{d}$ si et seulement si $ad - bc = 0$, et $\frac{a}{1} = a$. Alors $\text{Frac}(A)$ est un corps, l'addition et la multiplication étant définies de façon usuelle (comme dans $\mathbb{Q}$), et $A$ est un sous-anneau de $\text{Frac}(A)$.

Si $A \subseteq L$ est un sous-anneau d'un corps $L$, alors $A$ est intègre puisque $L$ l'est. Dans ce cas $\text{Frac}(A)$ est le plus petit sous-corps de $L$ contenant $A$ (exercice).

\begin{exemple}
$\text{Frac}(\mathbb{Z}) = \mathbb{Q}$ est le plus petit sous-corps de $\mathbb{C}$ (contenant $\mathbb{Z}$).
\end{exemple}

\begin{exemple}
Soit $\mathbb{F}_p[t]$ l'anneau des polynômes à coefficients dans $\mathbb{F}_p$ en la "variable" $t$. Son corps des fractions est
$$\text{Frac}(\mathbb{F}_p[t]) \triangleq \mathbb{F}_p(t) = \left\{\frac{P(t)}{Q(t)} \mid P(t), Q(t) \in \mathbb{F}_p[t], Q(t) \neq 0\right\}.$$
Les éléments de $\mathbb{F}_p(t)$ sont les "fractions rationnelles" à coefficients dans $\mathbb{F}_p$.
\end{exemple}

\textbf{Quotients.} Rappelons que si $I$ est un idéal de $A$, alors le quotient $A/I$ est un corps ssi $I$ est un idéal maximal. De plus, tout anneau commutatif non nul contient au moins un idéal maximal. Ainsi, si $\mathfrak{m}$ est un idéal maximal de $A$, alors $A/\mathfrak{m}$ est un corps.

\begin{exemple}
Les idéaux maximaux de $\mathbb{Z}$ sont les $p\mathbb{Z}$ avec $p$ premier. On obtient ainsi les corps finis à $p$ éléments $\mathbb{Z}/p\mathbb{Z} = \mathbb{F}_p$.
\end{exemple}

\begin{exemple}
Cet exemple est fondamental.

Soit $K$ un corps. L'anneau $K[X]$ est principal, et ses idéaux maximaux sont les $(P)$ avec $P(X)$ irréductible dans $K[X]$, que l'on peut choisir monique. Alors $K[X]/(P)$ est un corps, et aussi un $K$-espace vectoriel de dimension $\deg(P)$.

De plus, l'application $K \hookrightarrow K[X]/(P)$, $a \mapsto a \bmod (P)$ est un morphisme de corps injectif, qui identifie $K$ à un sous-corps de $K[X]/(P)$.
\end{exemple}

\section{Caractéristique d'un corps}
On distingue généralement deux types de corps. Soit $F$ un corps. Soit
$$\nu_F : \mathbb{Z} \longrightarrow F$$
$$1 \longmapsto 1_F$$
l'unique morphisme d'anneau de $\mathbb{Z}$ dans $F$. Par factorisation
$$\overline{\nu}_F : \mathbb{Z}/\ker(\nu_F) \hookrightarrow F$$
est un morphisme d'anneau injectif. Comme $F$ est un corps, il est intègre, et donc $\mathbb{Z}/\ker(\nu_F)$ aussi. Par conséquent $\ker(\nu_F)$ est un idéal premier de $\mathbb{Z}$, c'est-à-dire $\ker(\nu_F) = (0)$ ou $\ker(\nu_F) = p\mathbb{Z}$ avec $p$ premier.

$-$ Si $\ker(\nu_F) = p\mathbb{Z}$, alors on a un morphisme injectif
$$\overline{\nu}_F : \mathbb{Z}/p\mathbb{Z} = \mathbb{F}_p \hookrightarrow F$$
et $F$ contient un sous-corps isomorphe à $\mathbb{F}_p$.

$-$ Si $\ker(\nu_F) = (0)$, alors $\nu_F : \mathbb{Z} \hookrightarrow F$ s'étend en un morphisme injectif
$$\widetilde{\nu}_F : \mathbb{Q} \hookrightarrow F$$
$$\frac{a}{b} \longmapsto \nu_F(a)\nu_F(b)^{-1},$$
et $F$ contient un sous-corps isomorphe à $\mathbb{Q}$.

\begin{definition}
La caractéristique de $F$ est :

$-$ $\text{Car}(F) = p$ premier dans le cas où $\ker(\nu_F) = p\mathbb{Z}$ ;

$-$ $\text{Car}(F) = 0$ dans le cas où $\ker(\nu_F) = (0)$.
\end{definition}

\begin{remarque}
Si $\text{Car}(F) = p$ alors $p = 0$ dans $F$, et si $\text{Car}(F) = 0$ alors $n \neq 0$ dans $F$ pour tout $n \in \mathbb{Z} \setminus \{0\}$.
\end{remarque}

\begin{exemple}
$-$ Les corps $\mathbb{Q}$, $\mathbb{Q}(i)$, $\mathbb{R}$, $\mathbb{C}$ sont de caractéristique $0$. De façon générale, tous les sous-corps de $\mathbb{C}$ sont de caractéristique $0$.

$-$ Les corps $\mathbb{F}_p$, $\mathbb{F}_p(t)$ sont de caractéristique $p$.
\end{exemple}